% ============================================================================
% Section 6.4 — Constructing Triangles from Three Measurements
% CCSS 7.G.A.2
% ============================================================================
\section{Constructing Triangles from Three Measurements}

\begin{stepGoal}
\begin{itemize}
    \item Determine when three measurements give a unique triangle, many, or none
    \item Apply the triangle inequality and the $180°$ angle-sum rule
\end{itemize}
\end{stepGoal}

\begin{stepVocabBox}
    \vocabItem{SSS / SAS / ASA / AAS}{Shorthand for the given measurements --- S = side, A = angle.}
    \vocabItem{Angle Sum}{The three interior angles of every triangle add to $180°$.}
\end{stepVocabBox}

\begin{stepsCard}[How to Decide if Three Measurements Make a Triangle]
    \stepItem{Identify \textbf{what} is given: sides only (SSS), sides + angles (SAS, ASA, AAS), or angles only (AAA).}
    \stepItem{Check the \textbf{rules}: triangle inequality for sides; angle sum $= 180°$ for angles.}
    \stepItem{Use the table: SSS, SAS, ASA, AAS $\to$ \textbf{unique}; AAA $\to$ \textbf{many} (same shape, different sizes); rule fails $\to$ \textbf{none}.}
\end{stepsCard}

% --- Worked Example 1 ---
\begin{stepExample}{Can you draw a triangle with angles $50°$, $60°$, and $80°$?}
    \stepShow{1} All three measurements are angles (AAA).

    \stepShow{2} Angle sum: $50 + 60 + 80 = 190°$. This does not equal $180°$.

    \stepShow{3} The angle-sum rule fails.
    \stepResult{No triangle is possible.}
\end{stepExample}

% --- Worked Example 2 ---
\begin{stepExample}{Sides $7$, $10$, $15$. Can a triangle be formed? Is it unique?}
    \stepShow{1} Three sides given (SSS).

    \stepShow{2} Triangle inequality:\par
    $7 + 10 = 17 > 15$ \checkmark \quad $7 + 15 = 22 > 10$ \checkmark \quad $10 + 15 = 25 > 7$ \checkmark

    \stepShow{3} SSS with valid sides $\to$ exactly one triangle.
    \stepResult{Yes --- a unique triangle.}
\end{stepExample}

\mascotStepTip{If you know two angles of a triangle, you can always find the third: $180° - \text{angle 1} - \text{angle 2}$.}

% ============================================================================
% PRACTICE
% ============================================================================
\newpage
\begin{stepPractice}{Constructing Triangles Practice}
    \resetProblems

    \stepPracticeHeader[step1Color]{Identify the Type}
    \begin{multicols}{2}
    \prob Sides $3$, $4$, and angle $90°$ between them. Type? \answerBlank[2cm]
    \answer{SAS}
    \prob Angles $35°$, $65°$, $80°$. Type? \answerBlank[2cm]
    \answer{AAA}
    \end{multicols}

    \stepPracticeHeader[step2Color]{Check the Rules}
    \prob Do angles $35°$, $65°$, and $80°$ satisfy the angle-sum rule? \answerBlank[2cm]
    \answer{Yes ($35 + 65 + 80 = 180°$)}

    \prob Two angles of a triangle are $40°$ and $70°$. What is the third? \answerBlank[2cm]
    \answer{$70°$}

    \stepPracticeHeader[step3Color]{Unique, Many, or None?}
    \prob Sides $3$, $4$, included angle $90°$. Result? \answerBlank[2.5cm]
    \answer{Unique triangle (SAS)}

    \prob Angles $35°$, $65°$, $80°$. Result? \answerBlank[2.5cm]
    \answer{Many triangles (AAA --- same shape, different sizes)}

    \wordProblem{Can a triangle have angles $90°$, $90°$, and $0°$? Explain.}{~}
    \answer{No. Although $90 + 90 + 0 = 180°$, an angle of $0°$ collapses a side to nothing. Every angle must be greater than $0°$.}
\end{stepPractice}
